\subsubsection{Sinussatz}
\label{subsubsec:fundamentals_oblique_triangles_sine_rule}

Die Höhe $h$ zerlegt das schiefwinklige Dreieck in zwei rechtwinklige
Dreiecke. Dadurch lässt sich der Sinussatz aus den bereits bekannten
trigonometrischen Funktionen herleiten.

\begin{center}
  \CachedTikzFigure{content/resources/chapters/fundamentals/tikz/oblique_triangle_sine_rule_altitude.tex}
\end{center}

In den beiden rechtwinkligen Teildreiecken gilt
\begin{ClozeArea}[34]{9}
  \[
    \left.
    \begin{aligned}
      \text{linkes Dreieck:}\quad &\sin(\alpha)=\frac{h}{b}\\
      \text{rechtes Dreieck:}\quad &\sin(\beta)=\frac{h}{a}
    \end{aligned}
    \quad \right\}
    \quad\Longrightarrow\quad
    b\sin(\alpha)=h=a\sin(\beta)
    \quad\Longrightarrow\quad
    \frac{\sin(\alpha)}{a}=\frac{\sin(\beta)}{b}.
  \]
\end{ClozeArea}

Analog lässt sich dieser Zusammenhang für alle Seiten und Winkel des Dreiecks verallgemeinern.
\begin{TheoremBox}{Sinussatz}{fundamentals_sine_rule}
  In jedem Dreieck ist das Verhältnis des Sinus eines Winkels zur Länge der
  gegenüberliegenden Seite für alle drei Winkel gleich:
  \begin{ClozeArea}[20]{5}
    \[
      \frac{\sin(\alpha)}{a}=\frac{\sin(\beta)}{b}=\frac{\sin(\gamma)}{c}.
    \]
  \end{ClozeArea}
\end{TheoremBox}
